% Original PDF 21-22 / printed 11-12, complete Section 2.1.
% block: graded-chain-heading page:21 type:section
\subsection{{Graded vector spaces, chain complexes, and differentials}}\label{sec:graded-chain-complexes}

% block: graded-vector-space-definition page:21 type:ordinary-paragraph
{A \emph{graded vector space} is a vector space $V$ expressed as a direct sum of subspaces $V_k$ indexed by the integers:}

% block: graded-vector-space-sum page:21 type:equation
\[
V=\bigoplus_{k=-\infty}^{\infty}V_k.
\]

% block: graded-polynomial-spaces page:21 type:ordinary-paragraph
{For example, the space $\mathcal P(\mathbb R^n)$ of polynomial functions in $n$ variables is graded with the $k$th summand equal to the space of homogeneous polynomials of degree $k$ if $k\geq0$ and $0$ otherwise. A graded vector space is called \emph{nonnegative} if, as in this case, $V_k=0$ for $k<0$. Unlike the full space of polynomials, most of the graded spaces we encounter will be nonnegative and, in addition, \emph{finite}, meaning that only finitely many of the summands are nonzero. For example, the space $\mathcal P_r(\mathbb R^n)$ of polynomials of degree at most $r$ is graded with only $r+1$ nonzero summands.}

% block: graded-map-definition pages:21,22 type:ordinary-paragraph
{A linear map $f:V\to W$ between graded vector spaces is called a \emph{graded map of degree $p$} if $f(V_k)\subset W_{k+p}$ for $k$. The Laplace operator, for example, acts as a graded linear map of degree $-2$ on the polynomials, while multiplication by some fixed homogeneous polynomial of degree $p\geq0$ is a graded map of degree $p$.}

% block: chain-complex-definition page:22 type:ordinary-paragraph
{A \emph{chain complex} is simply a graded vector space $V$ furnished with a graded linear map $\partial:V\to V$ of degree $-1$ which satisfies $\partial^2=0$. The operator $\partial$, or, equivalently, the sequence of the operators $\partial_k=\partial|_{V_k}$, is called the \emph{differential} of the chain complex, or, sometimes, the \emph{boundary operator}. In other words, a chain complex consists of a sequence of vector spaces $V_k$ and linear maps $\partial_k:V_k\to V_{k-1}$, $k\in\mathbb Z$, with the property that $\partial_k\circ\partial_{k+1}=0$. We will often write a chain complex in the form}

% block: chain-complex-sequence page:22 type:equation
\[
\cdots\to V_{k+1}\xrightarrow{\partial_{k+1}}V_k\xrightarrow{\partial_k}V_{k-1}\to\cdots,
\]

% block: chain-complex-abbreviation page:22 type:ordinary-paragraph
{but when we wish to refer to it briefly, we will simply write $(V,\partial)$ or even just $V$. The elements of $V_k$ are sometimes called $k$-chains.}

% block: graded-subcomplex page:22 type:ordinary-paragraph
{A \emph{graded subspace} of a graded vector space is, naturally, a subspace $S$ which decomposes as the direct sum of subspaces $S_k\subset V_k$. If $S$ is a graded subspace of a complex $(V,\partial)$ such that $\partial S_k\subset S_{k-1}$, then $(S_k,\partial_k|_{S_k})$ is itself a complex, called a \emph{subcomplex} of $(V,\partial)$.}
